\documentclass[11pt,a4paper]{article}
\usepackage[utf8]{inputenc}
\usepackage[T1]{fontenc}
\usepackage[spanish]{babel}
\usepackage{geometry}
\usepackage{listings}
\usepackage{hyperref}
\usepackage{booktabs}
\usepackage{amsmath}

\geometry{margin=2.5cm}
\lstset{
    basicstyle=\ttfamily\small,
    breaklines=true,
    frame=single,
    numbers=left,
    numberstyle=\tiny,
    language=C++,
    keywordstyle=\bfseries,
    commentstyle=\itshape
}

\title{Práctico 1 --- Patrones de acceso a memoria\\
       \large Ejercicio 1: Localidad espacial}
\author{Pedro Calado Moura\\
        Documento: 65092945}
\date{}

\begin{document}
\maketitle

%==============================================================================
\section{Descripción del ejercicio}
%==============================================================================

El ejercicio consiste en implementar dos recorridas sobre un arreglo de gran tamaño (100 MB de \texttt{char}) que visitan exactamente la misma cantidad de elementos, pero con patrones de acceso diferentes:

\begin{enumerate}
\item \textbf{Recorrida secuencial:} se visita el arreglo en orden (posición 0, 1, 2, \ldots).
\item \textbf{Recorrida aleatoria:} se visitan las mismas posiciones pero en orden aleatorio.
\end{enumerate}

En ambos casos, el siguiente índice a visitar se obtiene leyendo de un arreglo de índices previamente inicializado. El objetivo es analizar el impacto de la \textbf{localidad espacial} en la jerarquía de memoria: el acceso secuencial explota la pre-carga de bloques contiguos en caché, mientras que el acceso aleatorio tiende a generar múltiples fallos de caché.

%==============================================================================
\section{Implementación}
%==============================================================================

Se utiliza un arreglo de 100 MB de \texttt{char} y dos arreglos de índices de tipo \texttt{uint32\_t} (suficiente para 100 millones de posiciones). Las mediciones de tiempo se realizan con \texttt{std::chrono::high\_resolution\_clock}.

\subsection{Estructura de datos y recorrida secuencial}

\begin{lstlisting}[caption={Reserva y recorrida secuencial}]
const size_t ARRAY_SIZE = 100 * 1024 * 1024;  // 100 MB
char* char_array = new char[ARRAY_SIZE];

std::vector<uint32_t> index_array(ARRAY_SIZE);
for (size_t i = 0; i < ARRAY_SIZE; ++i)
    index_array[i] = static_cast<uint32_t>(i);

auto start = std::chrono::high_resolution_clock::now();
for (size_t i = 0; i < ARRAY_SIZE; ++i) {
    uint32_t next_index = index_array[i];
    char_array[next_index] += 1;
}
auto end = std::chrono::high_resolution_clock::now();
\end{lstlisting}

\subsection{Recorrida con saltos aleatorios}

Los índices aleatorios se generan mediante una permutación usando el algoritmo Fisher-Yates (\texttt{std::shuffle}):

\begin{lstlisting}[caption={Indices aleatorios y recorrida}]
std::vector<uint32_t> random_index_array(ARRAY_SIZE);
for (size_t i = 0; i < ARRAY_SIZE; ++i)
    random_index_array[i] = static_cast<uint32_t>(i);
std::random_device rd;
std::mt19937 gen(rd());
std::shuffle(random_index_array.begin(), random_index_array.end(), gen);

auto start = std::chrono::high_resolution_clock::now();
for (size_t i = 0; i < ARRAY_SIZE; ++i) {
    uint32_t next_index = random_index_array[i];
    char_array[next_index] += 1;
}
\end{lstlisting}

%==============================================================================
\section{Resultados experimentales}
%==============================================================================

El programa se ejecutó en un sistema con carga reducida. Ambas recorridas visitan exactamente \(100 \times 1024 \times 1024 = 104\,857\,600\) elementos.

\begin{table}[h]
\centering
\begin{tabular}{@{}lc@{}}
\toprule
\textbf{Recorrida} & \textbf{Tiempo (ms)} \\
\midrule
Secuencial         & 69--80                \\
Aleatoria          & 1298--1317            \\
\bottomrule
\end{tabular}
\caption{Tiempos de ejecución en dos corridas.}
\label{tab:tiempos}
\end{table}

La recorrida aleatoria resulta aproximadamente \textbf{16--19 veces más lenta} que la secuencial. La variabilidad entre ejecuciones se debe a factores como el estado de la caché al inicio y la carga del sistema.

%==============================================================================
\section{Reflexión}
%==============================================================================

La diferencia de rendimiento se explica por la \textbf{localidad espacial} y la jerarquía de memoria.

\textbf{Acceso secuencial:} Cuando se accede a posiciones contiguas, la CPU aprovecha las \emph{cache lines} (bloques típicos de 64 bytes). Al leer una posición, se trae a L1/L2 todo el bloque contiguo; los accesos siguientes caen en caché sin necesidad de ir a memoria principal. El ancho de banda efectivo se maximiza.

\textbf{Acceso aleatorio:} Cada acceso puede dirigirse a una dirección distante que no está en caché. Esto provoca un \emph{cache miss} frecuente: la CPU debe ir a RAM para traer el dato. La memoria principal es del orden de 100 veces más lenta que la caché L1, por lo que el tiempo total aumenta drásticamente a pesar de realizar la misma cantidad de operaciones.

En resumen, el experimento demuestra que el \textbf{orden de acceso a los datos} afecta fuertemente el desempeño debido a la jerarquía de memoria. Los patrones de acceso que respetan la localidad espacial (como el secuencial) se benefician del caché y son significativamente más rápidos que los que dispersan los accesos por todo el espacio de direcciones.

\end{document}
